% =====================================================================
% LAMPIRAN A: GALERI TAMPILAN APLIKASI SMART CITY REPORTER
% =====================================================================
\chapter{Galeri Tampilan Aplikasi Smart City Reporter}
\label{app:ui-mockup}

Lampiran ini memuat galeri lengkap tampilan antarmuka untuk seluruh layar aplikasi \textit{Smart City Reporter} yang dirancang pada Bagian~\ref{sec:layar} Bab~\ref{ch:metode} dan dilaporkan hasil implementasinya pada Bagian~\ref{sec:hasil-flutter} Bab~\ref{ch:hasil}. Galeri disusun dari dua sumber tangkapan layar yang saling melengkapi. \emph{Pertama}, alur lengkap peran Warga Pelapor (Bagian~A.1--A.3) ditangkap langsung dari perangkat Android fisik saat aplikasi berjalan \textit{end-to-end} terhadap server inferensi FastAPI dan \textit{backend} Supabase, sehingga memuat data nyata, antara lain foto kejadian, koordinat dan alamat hasil \textit{reverse-geocoding}, nilai \textit{confidence}, instansi tujuan beserta estimasi jarak, dan metadata model. \emph{Kedua}, layar yang aksesnya dibatasi peran Operator Instansi dan Moderator atau \textit{Super Admin} (Bagian~A.4--A.7) dirender secara deterministik melalui \textit{golden screenshot test} (pengujian render \textit{widget} Flutter yang sebenarnya dengan data contoh terkendali), karena akun berperan instansi belum di-\textit{provisioning} pada perangkat uji dan data lintas-instansi tidak dapat ditampilkan dari satu akun warga. Pendekatan \textit{golden render} menjaga tampilan tetap konsisten dan dapat direproduksi tanpa bergantung pada koneksi jaringan. Pemisahan galeri ke lampiran menjaga pemaparan metodologi pada Bab~\ref{ch:metode} dan hasil pada Bab~\ref{ch:hasil} tetap padat.

% =====================================================================
\section*{A.1 Alur Pembuatan dan Klasifikasi Laporan (Perangkat Fisik)}

\begin{figure}[htbp]
    \centering
    \begin{subfigure}[b]{0.31\textwidth}
        \centering
        \includegraphics[height=6.2cm]{image-androdid/buat-laporan.jpg}
        \caption{Buat laporan}
    \end{subfigure}
    \hfill
    \begin{subfigure}[b]{0.31\textwidth}
        \centering
        \includegraphics[height=6.2cm]{image-androdid/Hasil-analisis-ai.jpg}
        \caption{Tinjauan hasil AI}
    \end{subfigure}
    \hfill
    \begin{subfigure}[b]{0.31\textwidth}
        \centering
        \includegraphics[height=6.2cm]{image-androdid/view-changes-instansi.jpg}
        \caption{Ganti instansi}
    \end{subfigure}
    \caption{Alur inti pelaporan oleh Warga Pelapor pada perangkat Android fisik: layar buat laporan (ambil foto kamera/galeri, narasi singkat, dan akuisisi lokasi GPS) (a); layar tinjauan hasil AI yang menampilkan instansi prediksi teratas beserta \textit{confidence} ($88\%$), peringkat alternatif, instansi tujuan terdekat dengan estimasi jarak, dan konfirmasi lokasi pada peta (b); serta lembar koreksi kategori manual ``Ganti instansi'' yang memungkinkan warga menimpa prediksi AI sebelum laporan dikirim (c).}
    \label{fig:app-pelaporan}
\end{figure}

% =====================================================================
\section*{A.2 Status, Riwayat, dan Detail Laporan (Perangkat Fisik)}

\begin{figure}[htbp]
    \centering
    \begin{subfigure}[b]{0.24\textwidth}
        \centering
        \includegraphics[height=5.6cm]{image-androdid/laporan-terkirim.jpg}
        \caption{Timeline status}
    \end{subfigure}
    \hfill
    \begin{subfigure}[b]{0.24\textwidth}
        \centering
        \includegraphics[height=5.6cm]{image-androdid/detail-laporan-terkirim.jpg}
        \caption{Detail teknis}
    \end{subfigure}
    \hfill
    \begin{subfigure}[b]{0.24\textwidth}
        \centering
        \includegraphics[height=5.6cm]{image-androdid/melihat-laporan.jpg}
        \caption{Riwayat laporan}
    \end{subfigure}
    \hfill
    \begin{subfigure}[b]{0.24\textwidth}
        \centering
        \includegraphics[height=5.6cm]{image-androdid/melihat-laporan-filter-terverifikasi.jpg}
        \caption{Filter status}
    \end{subfigure}
    \caption{Pelacakan laporan oleh warga: detail laporan dengan \textit{timeline} status \textit{submitted}~$\rightarrow$~\textit{verified}~$\rightarrow$~\textit{in progress}~$\rightarrow$~\textit{resolved} (a); panel ``Detail teknis'' yang memuat pelapor, prediksi dan \textit{confidence} AI, instansi tujuan, estimasi jarak \textit{assignment}, serta nama dan versi model (b); daftar riwayat ``Laporan Saya'' dengan kartu \textit{thumbnail}, kategori, dan status terkini (c); dan daftar yang sama setelah disaring ke status \textit{Terverifikasi} (d).}
    \label{fig:app-riwayat-detail}
\end{figure}

% =====================================================================
\section*{A.3 Profil Warga (Perangkat Fisik)}

\begin{figure}[htbp]
    \centering
    \includegraphics[height=7.4cm]{image-androdid/profile.jpg}
    \caption{Layar profil Warga Pelapor pada perangkat fisik: nama, surel, nomor telepon, dan peran (\textit{Citizen reporter}) yang diambil dari tabel \texttt{profiles} Supabase, beserta pintasan ke pengaturan dan keluar akun.}
    \label{fig:app-profil}
\end{figure}

% =====================================================================
\section*{A.4 Onboarding dan Autentikasi (Render Golden)}

\begin{figure}[htbp]
    \centering
    \begin{subfigure}[b]{0.235\textwidth}
        \centering
        \includegraphics[width=\linewidth]{bab_3/ui/ui_onboarding.png}
        \caption{Onboarding}
    \end{subfigure}
    \hfill
    \begin{subfigure}[b]{0.235\textwidth}
        \centering
        \includegraphics[width=\linewidth]{bab_3/ui/ui_login.png}
        \caption{Masuk akun}
    \end{subfigure}
    \hfill
    \begin{subfigure}[b]{0.235\textwidth}
        \centering
        \includegraphics[width=\linewidth]{bab_3/ui/ui_register.png}
        \caption{Registrasi}
    \end{subfigure}
    \hfill
    \begin{subfigure}[b]{0.235\textwidth}
        \centering
        \includegraphics[width=\linewidth]{bab_3/ui/ui_forgot.png}
        \caption{Lupa kata sandi}
    \end{subfigure}
    \caption{Layar onboarding tiga halaman serta layar masuk, registrasi akun terverifikasi, dan pemulihan kata sandi melalui surel.}
    \label{fig:app-onboarding}
\end{figure}

% =====================================================================
\section*{A.5 Beranda, Peta, dan Pengaturan (Render Golden)}

\begin{figure}[htbp]
    \centering
    \begin{subfigure}[b]{0.30\textwidth}
        \centering
        \includegraphics[width=\linewidth]{bab_3/ui/ui_dashboard.png}
        \caption{Beranda warga}
    \end{subfigure}
    \hfill
    \begin{subfigure}[b]{0.30\textwidth}
        \centering
        \includegraphics[width=\linewidth]{bab_3/ui/ui_map_view.png}
        \caption{Peta laporan}
    \end{subfigure}
    \hfill
    \begin{subfigure}[b]{0.30\textwidth}
        \centering
        \includegraphics[width=\linewidth]{bab_3/ui/ui_settings.png}
        \caption{Pengaturan}
    \end{subfigure}
    \caption{Beranda warga dengan ringkasan status laporan dan peta laporan terdekat; layar peta laporan dengan pin lokasi dan filter status di atas ubin OpenStreetMap; serta layar pengaturan aplikasi.}
    \label{fig:app-beranda-peta}
\end{figure}

% =====================================================================
\section*{A.6 Operator dan Admin Instansi (Render Golden)}

\begin{figure}[htbp]
    \centering
    \begin{subfigure}[b]{0.30\textwidth}
        \centering
        \includegraphics[width=\linewidth]{bab_3/ui/ui_agency_inbox.png}
        \caption{Antrean instansi}
    \end{subfigure}
    \hfill
    \begin{subfigure}[b]{0.30\textwidth}
        \centering
        \includegraphics[width=\linewidth]{bab_3/ui/ui_detail_admin.png}
        \caption{Perbarui status}
    \end{subfigure}
    \hfill
    \begin{subfigure}[b]{0.30\textwidth}
        \centering
        \includegraphics[width=\linewidth]{bab_3/ui/ui_evidence_sheet.png}
        \caption{Unggah bukti}
    \end{subfigure}
    \caption{Dasbor antrean Operator Instansi yang hanya menampilkan laporan ter-\textit{routing} sesuai peran melalui \textit{Row Level Security} (a); bilah aksi pembaruan status (\textit{verify}, \textit{reject}, \textit{in progress}, \textit{resolved}) saat operator membuka laporan instansinya (b); dan lembar unggah bukti penyelesaian yang wajib diisi sebelum laporan dapat ditandai selesai (c).}
    \label{fig:app-operator}
\end{figure}

\begin{figure}[htbp]
    \centering
    \begin{subfigure}[b]{0.30\textwidth}
        \centering
        \includegraphics[width=\linewidth]{bab_3/ui/ui_agency_admin_profile.png}
        \caption{Profil operator}
    \end{subfigure}
    \hspace{0.05\textwidth}
    \begin{subfigure}[b]{0.30\textwidth}
        \centering
        \includegraphics[width=\linewidth]{bab_3/ui/ui_detail_resolved.png}
        \caption{Laporan selesai}
    \end{subfigure}
    \caption{Profil operator yang memuat lingkup instansi yang ditugaskan (a), serta tampilan laporan yang telah selesai dengan \textit{timeline} status lengkap beserta bukti penyelesaian yang tersimpan (b).}
    \label{fig:app-operator-profil}
\end{figure}

% =====================================================================
\section*{A.7 Moderator atau Super Admin (Render Golden)}

\begin{figure}[htbp]
    \centering
    \begin{subfigure}[b]{0.30\textwidth}
        \centering
        \includegraphics[width=\linewidth]{bab_3/ui/ui_superadmin_inbox.png}
        \caption{Antrean lintas instansi}
    \end{subfigure}
    \hspace{0.05\textwidth}
    \begin{subfigure}[b]{0.30\textwidth}
        \centering
        \includegraphics[width=\linewidth]{bab_3/ui/ui_agency_registry.png}
        \caption{Registri instansi}
    \end{subfigure}
    \caption{Tampilan peran Moderator atau \textit{Super Admin}: antrean global ``All agencies'' yang menampilkan laporan dari seluruh instansi dengan tapis instansi dan panel \textit{triage} (a); serta layar registri instansi (CRUD atas tabel \texttt{agencies}) yang memuat nama kantor, kategori, koordinat, dan penanda aktif, lengkap dengan aksi pemuatan ulang server \textit{routing} (b). Akses kedua layar ini dibatasi peran \textit{super\_admin} melalui \textit{Row Level Security}.}
    \label{fig:app-superadmin}
\end{figure}
